\section{实验问题}
\subsection{sem\_wait阻塞进程}
在任务二实现进程互斥时在一开始使用函数\texttt{sem\_wait()}阻塞进程,该函数完全符合课上所说的P操作,即将信号量递减后判断是否小于零,如果是则阻塞当前进程,否则继续执行。
这样在实际运行文件系统时的确能正确实现读写进程间的互斥访问同一文件,但是一旦阻塞将导致当前交互界面无法接收执行任何命令,也就无法切换进程、释放资源等操作,系统正确但不可用。
\subsection{使用sem\_trywait替换}
考虑使用\texttt{sem\_trywait()}函数替换\texttt{sem\_wait()}函数,前者是后者的非阻塞尝试,成功获取信号量时操作相同,当失败时不会阻塞进程,而是立即返回,这样能实现任务二要求还能
保证系统的可用性,更符合此处实际场景。